% Genereller Header
\documentclass[10pt,twoside,a4paper,fleqn]{article}
\usepackage[left=1cm,right=1cm,top=1cm,bottom=1cm,includeheadfoot]{geometry}
\usepackage[ngerman]{babel,varioref}
\usepackage[utf8]{inputenc}

% Pakete
\usepackage{amssymb}
\usepackage{amsmath}
\usepackage{bm} % Bold Math
\usepackage{fancybox}
\usepackage{graphicx}
\usepackage[usenames,dvipsnames,svgnames,table]{xcolor}
\usepackage{lastpage}
\usepackage{wrapfig}
\usepackage{fancyhdr}
\usepackage{hyperref}
\usepackage{verbatim}
\usepackage{pdflscape} % landscape
\usepackage{multirow} % zellen in tabellen verbinden
\usepackage{multicol}
%\usepackage{slashbox} % getrennte zelle in tabelle
% \usepackage{array} % anordnung in tabellen
\usepackage{booktabs}
\usepackage{pgfplots}

%%%%%%%%%%%%%%%%%%%%
% Generelle Makros %
%%%%%%%%%%%%%%%%%%%%
\newcommand{\formelbuch}[1]{$_{\textcolor{red}{\mbox{\small{S#1}}}}$}
\newcommand{\verweis}[2]{ {\small (siehe auch \ref{#1}, #2 (S. \pageref{#1}))
}}
\newcommand{\subsubadd}[1]{\textcolor{black}{\mbox{#1}}}
\newenvironment{liste}[0]{
	\begin{list}{$\bullet$}{\setlength{\itemsep}{0cm}\setlength{\parsep}{0cm} \setlength{\topsep}{0cm}}}
    {\end{list}}
    
\newcommand{\logd}[0]{\log_{10}}
\newcommand{\subsubsubsection}[1]{\textbf{#1}}

\newenvironment{aufzaehlung}[0]{
	\begin{enumerate}{\setlength{\itemsep}{0cm}\setlength{\parsep}{0cm}
	\setlength{\topsep}{0cm}}} {\end{enumerate}}    

\newcommand{\abbHeight}[3]{
	\begin{center}
		\includegraphics[height=#2]{./bilder/#1} \\
		#3
    \end{center}
}
%Kreis um Buchstabe,Zahl
\newcommand*\kreisS[2]{\unitlength1ex\begin{picture}(2.5,2.5)%
\put(0.75,0.75){\circle{3}}\put(0.7,0.7){\makebox(0,0){#1}}\end{picture}$^{\displaystyle \text{#2}}$}

\newcommand*\kreisM[2]{\unitlength1ex\begin{picture}(3.5,3.5)%
\put(0.75,0.75){\circle{4}}\put(0.7,0.7){\makebox(0,0){#1}}\end{picture}$^{\displaystyle ^{\displaystyle \text{#2}}}$}

\newcommand*\kreisB[2]{\unitlength1ex\begin{picture}(3.5,3.5)%
\put(0.75,0.75){\circle{5}}\put(0.7,0.7){\makebox(0,0){#1}}\end{picture}$^{\displaystyle ^{\displaystyle \text{#2}}}$}

\newcommand{\todo}[1]{\colorbox{red}{#1}}
\newcommand{\hilight}[1]{\colorbox{yellow}{#1}}

%\newcommand{\skriptsection}[2]{\section{#1 {\tiny Skript S. #2}}}
%\newcommand{\skriptsubsection}[2]{\subsection{#1 {\tiny Skript S. #2}}}
%\newcommand{\skriptsubsubsection}[2]{\subsubsection{#1 {\tiny Skript S. #2}}}
%\renewcommand{\skriptsection}[2]{\section{#1 {\tiny Schaum S. #2}}}
%\renewcommand{\skriptsubsection}[2]{\subsection{#1 {\tiny Schaum S. #2}}}
%\renewcommand{\skriptsubsubsection}[2]{\subsubsection{#1 {\tiny Schaum S. #2}}}
\newcommand{\skriptsection}[2]{\section{#1 \formelbuch{#2}}}
\newcommand{\skriptsubsection}[2]{\subsection{#1 \formelbuch{#2}}}
\newcommand{\skriptsubsubsection}[2]{\subsubsection{#1 \formelbuch{#2}}}

%%%%%%%%%%
% Farben %
%%%%%%%%%%
\definecolor{black}{rgb}{0,0,0}
\definecolor{red}{rgb}{1,0,0}
\definecolor{white}{rgb}{1,1,1}
\definecolor{grey}{rgb}{0.8,0.8,0.8}

%%%%%%%%%%%%%%%%%%%%%%%%%%%%
% Mathematische Operatoren %
%%%%%%%%%%%%%%%%%%%%%%%%%%%%
\DeclareMathOperator{\sinc}{sinc}
\DeclareMathOperator{\sgn}{sgn}



% Fouriertransformationen
\unitlength1cm
\newcommand{\FT}
{
\begin{picture}(1,0.5)
\put(0.2,0.1){\circle{0.14}}\put(0.27,0.1){\line(1,0){0.5}}\put(0.77,0.1){\circle*{0.14}}
\end{picture}
}


\newcommand{\IFT}
{
\begin{picture}(1,0.5)
\put(0.2,0.1){\circle*{0.14}}\put(0.27,0.1){\line(1,0){0.45}}\put(0.77,0.1){\circle{0.14}}
\end{picture}
}



%%%%%%%%%%%%%%%%%%%%%%%%%%%%
% Allgemeine Einstellungen %
%%%%%%%%%%%%%%%%%%%%%%%%%%%%
%pdf info
\hypersetup{pdfauthor={\authorinfo},pdftitle={\titleinfo},colorlinks=false}
\author{\authorinfo}
\title{\titleinfo}

%Kopf- und Fusszeile
\pagestyle{fancy}
\fancyhf{}
%Linien oben und unten
\renewcommand{\headrulewidth}{0.5pt} 
\renewcommand{\footrulewidth}{0.5pt}


\fancyhead[L]{\titleinfo{ }- Summary}
%Kopfzeile rechts bzw. aussen
\fancyhead[R]{\today{ }- Page \thepage/\pageref{LastPage}}
\fancyfoot[L]{\includegraphics{bilder/ccbyncsa.png}}
\fancyfoot[R]{\authorinfo}

% Einr�cken verhindern versuchen
\setlength{\parindent}{0pt}

